\documentclass[../main]{subfiles}
\begin{document}
\ifSubfilesClassLoaded{\setcounter{page}{47}}{}
\section{Комунізм і особистість}
\label{sec:08_kommunizm_lichnost}

Проблема особистості – одна з центральних проблем комунізму. Значення особистості в умовах комунізму зростає до небувалих раніше висот. «Замість старого буржуазного суспільства з його класами та класовими протиріччями приходить асоціація, в якій вільний розвиток кожного є умовою вільного розвитку всіх», – писали Маркс і Енгельс у «Маніфесті комуністичної партії».

Усю попередню історію потрібно розуміти не просто як певну послідовність подій, а як процес формування сутності людини. Тоді стає зрозуміло, що це водночас процес формування певного типу особистості, яка відповідає сутності людини. Якою має бути ця особистість? Це питання щоразу вирішуватиметься практично і конкретно-історично, але в будь-якому випадку це особистість, яка свідомо творить історію в колективі таких же особистостей-творців. У вільній асоціації людей усе їхнє життя залежить від творчої діяльності, яка є найвищим проявом особистості, що її здійснює.

Особистість – явище, яке має соціальну природу від початку і до кінця. Особистістю людина не народжується, а стає, включаючись у суспільні відносини як активний учасник. Спочатку людське життя визначається для новонародженого ззовні – із соціального простору, в якому вона живе, з усіма варіаціями відносин людини до людини та людини до природи. І лише опанувавши, тобто зробивши своїми, форми діяльності або, що те саме, форми культури, які жодним чином спочатку не закладені в індивіді, вона стає особистістю. Ці форми кожен, хто з’являється на світ, знаходить вже готовими. Для біологічного організму вони є зовнішніми, але лише для організму, а не для особистості. Світ діяльності як світ людської культури – це справжній світ особистості. У комуністичному суспільстві кожен має зробити цей світ своїм, а це означає – навчитися змінювати його.

У комуністичному суспільстві особистість у своїй практиці активно створюватиме суспільні відносини, а не просто бути їхнім пасивним продуктом. Суть проблеми полягає в тому, щоб створити такі умови, в яких у своєму індивідуальному розвитку особистість якомога швидше пройшла б шлях становлення суспільства і вийшла на передову суспільного життя. Потрібно щоразу шукати та знаходити конкретні форми колективності, які б пов’язували особистість із суспільством таким чином, щоб вона не могла не бути його активним членом і не жила всім багатством суспільного життя. Оскільки суспільство постійно розвивається, такі форми «спільної окремої діяльності» не можуть бути знайдені один раз і назавжди. Управління суспільним розвитком за комунізму полягає, головним чином, у тому, щоб навчитися вчасно змінювати форми зв’язку людини з суспільством таким чином, щоб вони змушували її діяти по-людськи в будь-яких нових конкретно-історичних умовах.

Комуністів звинувачують у тому, що вони прагнуть стандартизувати особистість, зробити людей «однаковими». У цьому звинуваченні прихильники капіталізму, по суті, розкривають таємницю сучасного капіталістичного суспільства. За часів капіталізму індивіди в своїй масі настільки стандартизовані, що їхня індивідуальність визначається якимись дрібницями. Ця хибна індивідуальність не має жодного стосунку до справжньої індивідуальності, як до індивідуального вираження засвоєної людиною загальної людської культури. Капіталізм залучає людей до стандартних практик виробництва та споживання, стандартизує їхню поведінку як під час роботи, так і поза нею. Люди вже «однакові». Комунізм же передбачає створення суспільного простору для розвитку індивідуальності.

Комуністи вважають необхідним створення суспільства, в якому розвиватимуться всебічно обдаровані, творчі особистості. У такому суспільстві індивідуальність кожної людини буде активним вираженням і втіленням суспільного розвитку.

В якому сенсі тут можна говорити про стандарти? Очевидно, людина не може бути всебічно розвиненим творцем, не опанувавши необхідний для цього мінімум людської культури, і, відповідно, слід говорити про мінімум умов, які суспільство має забезпечити кожному своєму членові, щоб він міг бути всебічно розвиненим творцем. Який саме це мінімум?

До цього належать: певний рівень розвитку системи охорони здоров’я, і саме спрямованої на збереження здоров’я та створення умов для здорового життя кожної людини, а не лише на лікування хвороб; якість і рівень освіти, обов’язкові для кожного; система громадського харчування та загалом усіх загальнодоступних служб, які можуть покращити побут; забезпечення дитинства та старості. В одній зі своїх праць Ленін писав, що соціалізм починається зі склянки молока для кожної дитини. Очевидно, що наявність цієї склянки ще не свідчить про те, що суспільство розвивається в напрямку комунізму, але без того, щоб діти були ситими, воно точно не може розвиватися в цьому напрямку. Так само, як і не може розвиватися без усього вищезазначеного (цей перелік, звісно, не повний. Його розробкою повинні займатися люди, які розбираються в різних галузях, і в міру розвитку суспільства він буде розширюватися).

Цей мінімум людяного ставлення до людини не може бути остаточно встановлений. Він змінюється залежно від умов, у яких відбувається перехід: від рівня розвитку людської культури. У сучасний мінімум, який став би поштовхом для розвитку суспільства до комунізму, обов’язково має входити вища освіта, що включає не лише комп’ютерну грамотність на рівні програмування, але й наукові основи розуміння всіх виробничих процесів – політехнізм (для всіх!). Також обов’язковим є забезпечення доступу до всіх сучасних засобів комунікації.

Важливо зазначити, що навіть найбільш розвинені капіталістичні країни не можуть похвалитися повною грамотністю всього населення та доступністю медичного обслуговування.

Необхідний мінімум людяності у ставленні до людини, як вже зазначалося, є явищем, що має свою історичну еволюцію. Але цей мінімум також має включати необхідність його постійного розширення. Щоб комуністичне суспільство, де вільний розвиток кожної особи є умовою розвитку всіх, могло існувати та розвиватися, в цей мінімум повинні постійно включатися всі нові передові досягнення людства. Комуністичне суспільство складатиметься з розумних, здорових і талановитих людей, які, зрештою, зможуть організувати суспільне відтворення свого життя на розумних засадах.
\end{document}
